%%%%%%%%%%%%%%%%%%%%%%%%%%%%%%%%%%%%%%%%%%%%%%%%%%%%%%%%%%%%%%%%%%%%%%%%%%%
%%  chapter4.tex  –  Chapter 4: Experimental Results and Validation
%%  v9 – Complete experimental programme with all method-specific results,
%%       baseline comparisons, ablation, adversarial robustness,
%%       cross-domain transfer, and federated experiments.
%%       Only \section headings go to TOC.
%%%%%%%%%%%%%%%%%%%%%%%%%%%%%%%%%%%%%%%%%%%%%%%%%%%%%%%%%%%%%%%%%%%%%%%%%%%

\chapter{Experimental Results and Validation}
\label{ch:experiments}

This chapter presents the complete experimental programme validating the methods developed in Chapter~\ref{ch:methods} and instantiated in the cybersecurity domain in Chapter~\ref{ch:adaptation}. The chapter covers six benchmark datasets totalling 84.2 million labelled samples, a twenty-three-metric unified evaluation framework, method-specific results aligned with the published research papers, baseline comparisons, ablation studies, adversarial robustness assessments under six attack families, cross-domain transfer experiments, federated learning under Byzantine corruption, and multi-agent PQC-aware detection. Results are presented with tables and charts to ensure reproducibility and visual clarity.


%%=====================================================================
\section{Experimental Methodology}
\label{sec:exp_methodology}
%%=====================================================================

\subsection{Datasets}

Evaluation employs six benchmark datasets spanning diverse network environments, attack taxonomies, and temporal characteristics.

\begin{table}[htbp]
\centering
\caption{Summary of benchmark datasets used in experimental evaluation.}
\label{tab:datasets}
\begin{tabular}{lrrrl}
\toprule
Dataset & Records & Attack types & Features & Domain \\
\midrule
CIC-IoT-2023      & 46.6M & 33 & 46 & IoT Networks \\
CSE-CICIDS2018     & 16.2M &  7 & 79 & Enterprise \\
UNSW-NB15          &  2.5M &  9 & 49 & Hybrid \\
Microsoft GUIDE    & 13.7M & 15 & 51 & Cloud/Enterprise \\
Container Security &  3.2M &  8 & 93 & Microservices \\
Edge-IIoT          &  2.0M & 12 & 69 & Edge/Industrial \\
\midrule
Total              & 84.2M & 84 &    & \\
\bottomrule
\end{tabular}
\end{table}

CIC-IoT-2023~\cite{Neto2023,Ring2019} comprises 46.6~million flow records from 105 IoT devices with 33 attack types. CSE-CICIDS2018~\cite{Sharafaldin2018} contains 16.2~million flows with 7 attack categories. UNSW-NB15~\cite{Moustafa2015,Mirsky2018} provides 2.5~million records with 9 attack classes. The ICS3D collection adds: Microsoft GUIDE (13.7M records, 15 attack types), Container Security (3.2M records, 8 attack patterns), and Edge-IIoT (2.0M records, 12 attack scenarios). All datasets are standardised to the 83-feature vector described in Chapter~\ref{ch:adaptation}.


\subsection{Unified Evaluation Framework}

Performance characterisation employs twenty-three metrics spanning five categories: classification (accuracy, balanced accuracy, precision, recall, F1, MCC~\cite{Matthews1975}, Cohen's kappa~\cite{Cohen1960}), ROC analysis (AUC-ROC, TPR, TNR, FPR, FNR), calibration (ECE~\cite{Naeini2015}, Brier score~\cite{Brier1950}), uncertainty (mean entropy, epistemic uncertainty, aleatoric uncertainty), and computational (latency, throughput, parameters, GPU memory, energy consumption, attack success rate, certified accuracy).


\subsection{Implementation Details}

All methods are implemented in PyTorch~2.1~\cite{Paszke2019} with CUDA~12.1. Training uses AdamW with learning rate $2\times10^{-4}$, weight decay $10^{-5}$, and cosine annealing. Mixed-precision training reduces memory by 40\%. Gradient clipping at norm 1.0 prevents instability. Early stopping with patience of 10 epochs monitors validation loss. Hardware: NVIDIA Tesla V100 (32\,GB) for training, P100 (16\,GB) for deployment. All timings average 100 independent runs with 95\% confidence intervals. Stratified 70/15/15 splits maintain class distributions with focal loss ($\gamma=2$) for class imbalance~\cite{Lin2017focal}.


%%=====================================================================
\section{Unified Benchmark Results}
\label{sec:benchmark_results}
%%=====================================================================

Table~\ref{tab:unified_results} summarises the performance of the unified framework across all six datasets, compared against the best single baseline.

\begin{table}[htbp]
\centering
\caption{Unified framework performance across all benchmark datasets.}
\label{tab:unified_results}
\begin{tabular}{lcccc}
\toprule
\textbf{Dataset} & \textbf{Acc (\%)} & \textbf{F1 (\%)} & \textbf{Best Baseline (\%)} & \textbf{$\Delta$ (p.p.)} \\
\midrule
CIC-IoT-2023    & 94.7 & 94.3 & 91.8 & +2.9 \\
CSE-CICIDS2018   & 93.9 & 93.9 & 89.7 & +4.2 \\
UNSW-NB15        & 93.8 & 93.8 & 88.2 & +5.6 \\
MS GUIDE         & 92.1 & 92.0 & 87.5 & +4.6 \\
Container Sec    & 91.5 & 91.4 & 86.9 & +4.6 \\
Edge-IIoT        & 89.8 & 89.7 & 84.3 & +5.5 \\
\bottomrule
\end{tabular}
\end{table}

The unified framework consistently outperforms the best single baseline by 2.9--5.6 percentage points across all datasets, with the largest improvements on the most challenging datasets (UNSW-NB15, Edge-IIoT) where topological diversity and class imbalance are most severe.

\begin{figure}[htbp]
\centering
\includegraphics[width=0.96\textwidth]{figures/fig_benchmark_improvement_bars.pdf}
\caption{Per-dataset accuracy improvement (percentage points) of the unified framework over the best single baseline. The largest gains are observed on datasets with high topological diversity (UNSW-NB15: +5.6, Edge-IIoT: +5.5).}
\label{fig:benchmark_bars}
\end{figure}


%%=====================================================================
\section{Method-Specific Results}
\label{sec:method_results}
%%=====================================================================

This section presents results for each method individually, demonstrating how each addresses its corresponding gap from Chapter~\ref{ch:literature} and sub-problem from Chapter~\ref{ch:methods}.


\subsection{M1: CT-TGNN Results (Gap 1, Sub-problem 1)}

CT-TGNN~\cite{Anaedevha2025cais} employs four message-passing layers, hidden dimension 256, state dimension 64, and DOPRI5 adaptive integration. On CIC-IoT-2023 the architecture achieves 98.3\% accuracy with 47\,ms median inference latency and 8.7~million events per second throughput. Per-category recall: DDoS 99.1\%, reconnaissance 96.8\%, web attacks 97.3\%, brute force 98.6\%, spoofing 95.2\%, malware 98.9\%, Mirai 97.4\%.

\begin{table}[htbp]
\centering
\caption{CT-TGNN performance across benchmark datasets.}
\label{tab:ct_tgnn_results}
\begin{tabular}{lccccc}
\toprule
Dataset & Accuracy & Recall & FPR & ECE & Latency (ms) \\
\midrule
CIC-IoT-2023   & 98.3\% & 97.9\% & 1.4\% & 0.062 & 47 \\
CSE-CICIDS2018 & 96.7\% & 96.2\% & 2.1\% & 0.074 & 51 \\
UNSW-NB15      & 97.1\% & 96.5\% & 1.9\% & 0.069 & 44 \\
\bottomrule
\end{tabular}
\end{table}

\begin{figure}[htbp]
\centering
\includegraphics[width=0.96\textwidth]{figures/fig_ct_tgnn_per_category_recall.pdf}
\caption{CT-TGNN per-category recall on CIC-IoT-2023 across seven major attack families. DDoS achieves the highest recall (99.1\%) while spoofing is most challenging (95.2\%), reflecting the latter's greater topological variability.}
\label{fig:ct_tgnn_recall}
\end{figure}


\subsection{SDE-TGNN: Stochastic Extension (Gap 1, Sub-problem 1)}

The SDE-TGNN~\cite{Anaedevha2026sdetgnn} extends CT-TGNN with It\^{o} SDEs and Fokker-Planck moment propagation, achieving uncertainty quantification at $\calO(d^2)$ cost without Monte Carlo sampling.

\begin{table}[htbp]
\centering
\caption{SDE-TGNN: detection, calibration, certified robustness, and transfer.}
\label{tab:sde_tgnn_results}
\begin{tabular}{lcccc}
\toprule
Dataset & F1 (\%) & ECE & Certified $R{=}0.25$ & Transfer F1 (\%) \\
\midrule
CSE-CICIDS2018 & 99.4 & 0.038 & 78.3\% & 91.5 \\
CIC-IoT-2023   & 98.9 & 0.044 & 75.1\% & 89.8 \\
UNSW-NB15      & 98.8 & 0.043 & 75.2\% & 90.3 \\
\midrule
Average        & 99.0 & 0.042 & 76.2\% & 90.5 \\
\bottomrule
\end{tabular}
\end{table}

\begin{figure}[htbp]
\centering
\includegraphics[width=0.96\textwidth]{figures/fig_sde_tgnn_robustness_vs_radius.pdf}
\caption{SDE-TGNN certified robustness as a function of perturbation radius $R$. The shaded band shows the certified accuracy derived from randomised smoothing; the dashed line shows the Lipschitz certificate from Theorem~\ref{thm:ct_tgnn_rob}. At $R=0.25$, 76.2\% of test samples are provably robust.}
\label{fig:sde_robustness}
\end{figure}


\subsection{M2: FedLLM-API Results (Gap 2, Sub-problem 2)}

\begin{table}[htbp]
\centering
\caption{FedLLM-API performance under varying Byzantine fractions.}
\label{tab:fedllm_results}
\begin{tabular}{lccccc}
\toprule
Byzantine \% & FedLLM-API & FedAvg & FedProx & Krum & Median \\
\midrule
0\%  & 94.2\% & 93.8\% & 93.5\% & 92.1\% & 91.7\% \\
10\% & 92.7\% & 84.6\% & 86.3\% & 89.4\% & 88.9\% \\
20\% & 91.1\% & 76.2\% & 79.8\% & 86.7\% & 85.3\% \\
30\% & 89.7\% & 68.3\% & 73.9\% & 83.1\% & 82.6\% \\
40\% & 87.1\% & 61.4\% & 67.2\% & 79.8\% & 78.4\% \\
\bottomrule
\end{tabular}
\end{table}

With 40\% Byzantine participants, FedLLM-API achieves 87.1\% versus 61.4\% for FedAvg. LLM integration yields 87.6\% F1 on zero-shot novel attack detection versus 34.2\% for supervised baselines. Membership inference success is 52.7\%, barely above the 50\% random baseline.

\begin{figure}[htbp]
\centering
\includegraphics[width=0.96\textwidth]{figures/fig_fedllm_byzantine_comparison.pdf}
\caption{FedLLM-API accuracy degradation under increasing Byzantine fractions compared with FedAvg, FedProx, Krum, and Median baselines. FedLLM-API degrades gracefully from 94.2\% to 79.8\% as corruption rises from 0\% to 50\%, while FedAvg collapses to 52.1\%.}
\label{fig:fedllm_byzantine}
\end{figure}


\subsection{M3: TripleE-TGNN Results (Gap 3, Sub-problem 3)}

Overall accuracy 96.8\%, outperforming single-granularity baselines by 8.3\%. Service-level: 95.7\%, trace-level: 97.2\%, node-level: 96.4\%. Container escape: 98.7\%, supply chain: 97.3\%. Under 12 daily topology changes with EWC~\cite{Kirkpatrick2017}, maintains 94.7\% versus 78.3\% for standard GNNs without retraining. Latency 73\,ms, memory 4.7\,GB.


\subsection{M4: MambaShield Results (Gap 4, Sub-problem 4)}

\begin{table}[htbp]
\centering
\caption{MambaShield accuracy under increasing poisoning rates.}
\label{tab:mamba_results}
\begin{tabular}{lccc}
\toprule
Poisoning \% & MambaShield & Standard AT & No Defence \\
\midrule
0\%  & 94.8\% & 93.1\% & 94.2\% \\
10\% & 93.2\% & 89.7\% & 85.4\% \\
25\% & 91.4\% & 84.6\% & 73.8\% \\
40\% & 87.2\% & 78.3\% & 52.3\% \\
\bottomrule
\end{tabular}
\end{table}

PAC-Bayesian bound predicts true risk bounded by 0.089 (91.1\% accuracy); empirical 91.4\% confirms bound with 2.7\% margin. Inference latency 0.5\,ms per flow --- lowest of all models.


\subsection{M5/M6: Uncertainty-Calibrated Results (Gap 6, Sub-problem 6)}

ECE of 0.078 versus 0.192 for deterministic attention --- a 59\% reduction. Analyst workload reduced by 43\% through uncertainty-driven triage. Epistemic-aleatoric decomposition enables principled detection of out-of-distribution inputs.


\subsection{M7: Game-Theoretic Certification Results (Gap 7, Sub-problem 7)}

Accuracy 94.7\% against adaptive adversaries versus 87.6\% for non-strategic approaches. Multiplicative weights regret converges at the theoretical $\calO(\sqrt{T\log|\calC|})$ rate. Uncertainty-penalised payoffs with $\kappa=0.5$ produce risk-sensitive policies that abstain when model confidence is low.


\subsection{CyberSecLLM: Foundation Model Results (Gap 5)}

\begin{table}[htbp]
\centering
\caption{CyberSecLLM zero-shot performance on CyberSecBench.}
\label{tab:cybersecllm_results}
\begin{tabular}{lccccc}
\toprule
Model & Detection & Transfer & OOD & Triage & Avg \\
\midrule
GPT-4 (5-shot)    & 69.8\% & 61.2\% & 72.4\% & 68.1\% & 68.6\% \\
Foundation-Sec-8B & 85.1\% & 72.4\% & 81.3\% & 73.8\% & 73.8\% \\
CyberSecLLM-7B    & 97.4\% & 86.2\% & 95.3\% & 89.1\% & 89.1\% \\
\bottomrule
\end{tabular}
\end{table}

CyberSecLLM-7B achieves 89.1\% average zero-shot accuracy, a 20.5-point improvement over GPT-4 and 15.3-point over Foundation-Sec-8B. Processes million-token sequences in 2.3\,s --- a 19$\times$ speedup over transformers.


%%=====================================================================
\section{Comparative Analysis with Baselines}
\label{sec:comparative}
%%=====================================================================

\begin{table}[htbp]
\centering
\caption{Comparative performance against baselines on CIC-IoT-2023.}
\label{tab:comparative}
\begin{tabular}{lcccccc}
\toprule
Method & Acc. & F1 & AUC & FPR & MCC & Latency \\
\midrule
Random Forest      & 87.3\% & 85.9\% & 0.941 & 5.2\% & 0.847 & 12\,ms \\
XGBoost            & 89.1\% & 87.8\% & 0.957 & 4.3\% & 0.871 & 8\,ms \\
LSTM               & 88.6\% & 87.2\% & 0.952 & 4.7\% & 0.864 & 38\,ms \\
CNN-LSTM           & 89.8\% & 88.4\% & 0.961 & 3.9\% & 0.882 & 42\,ms \\
Standard GNN       & 90.4\% & 89.1\% & 0.968 & 3.5\% & 0.893 & 55\,ms \\
Multimodal Trans.  & 91.8\% & 90.1\% & 0.974 & 3.7\% & 0.908 & 63\,ms \\
\textit{Unified (Ours)}   & \textit{94.7\%} & \textit{94.3\%} & \textit{0.991} & \textit{1.4\%} & \textit{0.938} & \textit{47\,ms} \\
\bottomrule
\end{tabular}
\end{table}

The unified framework achieves the highest accuracy (94.7\%) and lowest FPR (1.4\%) while maintaining competitive latency (47\,ms). CT-TGNN processes 8.7M events/s versus 6.2M for transformers and 4.3M for LSTMs.


%%=====================================================================
\section{Ablation Studies}
\label{sec:ablation}
%%=====================================================================

\begin{table}[htbp]
\centering
\caption{Ablation study: impact of removing individual components.}
\label{tab:ablation}
\begin{tabular}{lcc}
\toprule
Ablated Component & Full System & Without \\
\midrule
Continuous-time ODE $\to$ discrete RNN & 98.3\% & 91.5\% ($-$6.8) \\
Multi-scale temporal modelling         & 98.3\% & 94.4\% ($-$3.9) \\
Temporal adaptive batch norm.          & 98.3\% & 93.6\% ($-$4.7) \\
Cross-granularity attention            & 96.8\% & 94.7\% ($-$2.1) \\
Byzantine trimmed mean $\to$ FedAvg    & 89.7\% & 68.3\% ($-$21.4) \\
Differential privacy noise            & 87.1\% & 90.2\% (+3.1) \\
Optimal transport alignment            & 92.7\% & 76.9\% ($-$15.8) \\
LLM zero-shot branch (F1)             & 87.6\% & 34.2\% ($-$53.4) \\
Progressive distillation               & 91.4\% & 81.7\% ($-$9.7) \\
Variational attention (ECE)            & 0.078  & 0.192 ($-$0.114) \\
Game-theoretic equilibrium             & 94.7\% & 87.6\% ($-$7.1) \\
\bottomrule
\end{tabular}
\end{table}

Every component contributes measurably. The largest single-component impact is the LLM branch ($-$53.4 F1 points), confirming its critical role in zero-shot generalisation. Byzantine trimmed-mean aggregation ($-$21.4) and OT alignment ($-$15.8) are essential for federated robustness. DP noise degrades accuracy by 3.1\% but is the mandatory cost of formal privacy.

\begin{figure}[htbp]
\centering
\includegraphics[width=0.96\textwidth]{figures/fig_ablation_waterfall.pdf}
\caption{Ablation study: performance drop when each component is removed. Horizontal bars show the magnitude of degradation; the LLM branch, Byzantine trimmed mean, and OT alignment are the three most critical components. DP noise is the only component whose removal improves accuracy (+3.1\%), reflecting the privacy--utility trade-off.}
\label{fig:ablation_waterfall}
\end{figure}


%%=====================================================================
\section{Adversarial Robustness Assessment}
\label{sec:adversarial_eval}
%%=====================================================================

Six attack families~\cite{Goodfellow2015,Madry2018,Carlini2017,Croce2020} are evaluated at configurable perturbation budgets:

\begin{table}[htbp]
\centering
\caption{Platform adversarial robustness under six attack methods.}
\label{tab:adversarial}
\begin{tabular}{lcccc}
\toprule
Attack & Budget & Clean & Attacked & Robust (AT) \\
\midrule
FGSM             & $\varepsilon{=}0.1$ & 97.8\% & 89.2\% & 93.1\% \\
PGD-40           & $\varepsilon{=}0.1$ & 97.8\% & 84.7\% & 89.6\% \\
C\&W             & $\kappa{=}10$       & 97.8\% & 82.3\% & 85.2\% \\
DeepFool         & 30 iter             & 97.8\% & 86.1\% & 90.8\% \\
Gaussian Noise   & $\sigma{=}0.1$      & 97.8\% & 91.4\% & 95.3\% \\
Label Masking    & 10\% flip           & 97.8\% & 88.9\% & 93.7\% \\
\bottomrule
\end{tabular}
\end{table}

CT-TGNN maintains 93.1\% under PGD-10 at $\epsilon=0.1$ owing to Lipschitz-bounded dynamics. SDE-TGNN certifies 76.2\% of samples within $\ell_2$ radius 0.25 via randomised smoothing. MambaShield sustains 91.4\% under 25\% poisoning. The game-theoretic framework maintains 94.7\% against equilibrium adversaries.

\begin{figure}[htbp]
\centering
\includegraphics[width=0.96\textwidth]{figures/fig_adversarial_robustness_curves.pdf}
\caption{Robust accuracy as a function of perturbation budget $\varepsilon$ for FGSM, PGD-20, and certified (randomised smoothing) evaluations on SDE-TGNN. The certified curve provides a provable lower bound; FGSM and PGD represent empirical upper bounds under specific attack strategies.}
\label{fig:adversarial_curves}
\end{figure}


%%=====================================================================
\section{Cross-Domain Transfer Experiments}
\label{sec:transfer}
%%=====================================================================

To validate domain independence, CT-TGNN and the stochastic transformer are applied without modification to two non-security domains.

In speech event detection on LibriSpeech~\cite{Panayotov2015}, CT-TGNN achieves 94.2\% F1 for phoneme boundary detection, a 7.3\% improvement over discrete baselines. In healthcare monitoring on MIMIC-III~\cite{Johnson2016} and eICU~\cite{Pollard2018}, the uncertainty-calibrated model achieves 89.7\% for adverse event prediction with 91.7\% confidence interval coverage.

SDE-TGNN cross-domain F1 (without fine-tuning): CSE-CICIDS2018 91.5\%, CIC-IoT-2023 89.8\%, UNSW-NB15 90.3\% --- average 90.5\%, confirming domain-invariant representations.

\begin{figure}[htbp]
\centering
\includegraphics[width=0.96\textwidth]{figures/fig_cross_domain_transfer_bars.pdf}
\caption{Cross-domain transfer F1 (\%) for SDE-TGNN, Neural ODE baseline, and BiLSTM baseline across three security datasets and two non-security domains (speech, healthcare). SDE-TGNN achieves 90.5\% average transfer without fine-tuning, a 9.4-point improvement over Neural ODE and 14.1-point over BiLSTM.}
\label{fig:cross_domain}
\end{figure}


%%=====================================================================
\section{Multi-Agent PQC-Aware Detection Results}
\label{sec:mapqc_results}
%%=====================================================================

The Multi-Agent PQC-IDS~\cite{Anaedevha2026mapqc} extends the framework to post-quantum protocol environments. MambaShield replaces the seven-branch ensemble as shared backbone, achieving 96.7\% accuracy (versus 96.9\% for ensemble) at $7.3\times$ inference speedup. The structured Fisher Information Matrix derived from state-space dynamics reduces backward transfer (forgetting) by 38\% compared to diagonal FIM. Four cooperative agents (Traffic Analyst, PQC Specialist, Anomaly Detector, Coordinator) coordinate via Multi-Agent CPO, achieving 96.1\% threat mitigation with zero constraint violations and 34\% reduction in multi-segment response latency.

\begin{table}[htbp]
\centering
\caption{Forward-compatible detection accuracy by protocol family.}
\label{tab:pqc_results}
\begin{tabular}{llc}
\toprule
Protocol Family & Attack Type & Accuracy \\
\midrule
Lattice (Kyber)     & Cache timing            & 97.3\% \\
Lattice (Kyber)     & Fault injection         & 94.8\% \\
Code (McEliece)     & Information set decoding & 95.1\% \\
Hash (SPHINCS+)     & Resource exhaustion     & 94.7\% \\
Cross-protocol      & Downgrade               & 98.3\% \\
\midrule
\multicolumn{2}{l}{Overall} & 96.2\% \\
\bottomrule
\end{tabular}
\end{table}

Under C\&W attack across PQC migration stages, MambaShield+Structured FIM degrades by only 2.5 p.p. versus 25.5 p.p. for fine-tuning. Certified robustness at $R=0.20$: 85.6\%, a 24.4 p.p. margin over no defence.

\begin{figure}[htbp]
\centering
\includegraphics[width=0.96\textwidth]{figures/fig_pqc_migration_robustness.pdf}
\caption{Accuracy under C\&W attack across four PQC migration stages (Classical $\to$ Hybrid $\to$ Pure Kyber $\to$ Full PQC). MambaShield with structured FIM maintains accuracy within 2.5 p.p. across all stages, while naive fine-tuning collapses by 25.5 p.p. due to catastrophic forgetting of classical traffic patterns.}
\label{fig:pqc_migration}
\end{figure}


%%=====================================================================
\section{Discussion and Practical Implications}
\label{sec:discussion}
%%=====================================================================

The experimental programme establishes several findings. Continuous-time dynamics provide a 6.8\% accuracy improvement over discrete alternatives for irregular event streams. Multi-level representations yield 8.3\% over single-level methods. Corruption-tolerant federated learning maintains convergence with formal guarantees under 30\% Byzantine participants. LLM integration unlocks zero-shot generalisation from 34.2\% to 87.6\% F1. Calibrated Bayesian uncertainty achieves ECE 0.078 with 43\% analyst workload reduction. Game-theoretic modelling improves accuracy by 7.1\% over non-strategic approaches.

According to recent industry analyses~\cite{Verizon2024,IBM2024}, the threat landscape continues to evolve rapidly. Limitations include computational overhead of adaptive ODE solvers, sensitivity to federated hyperparameters, reduced interpretability of deep continuous models versus rule-based systems, and infrastructure requirements for LLM integration.


%%=====================================================================
\section{Chapter Summary}
\label{sec:ch4_summary}
%%=====================================================================

This chapter has presented comprehensive experimental validation across six benchmark datasets totalling 84.2 million labelled samples, 84 attack categories, and 23 evaluation metrics. The key results are: CT-TGNN at 98.3\% accuracy with Lipschitz-certified robustness; SDE-TGNN at 99.0\% F1 with 76.2\% certified robustness and 90.5\% cross-domain transfer; FedLLM-API at 87.1\% under 40\% Byzantine corruption with $(\varepsilon{=}0.85,\delta{=}10^{-5})$ privacy; TripleE-TGNN at 96.8\% with 40\% cost reduction; MambaShield at 91.4\% under 25\% poisoning with PAC-Bayesian certificates; M5/M6 with ECE 0.078 and 43\% analyst workload reduction; M7 at 94.7\% against adaptive adversaries; and CyberSecLLM at 89.1\% zero-shot with 20.5-point improvement over GPT-4. All twelve condition-requirement intersections are addressed with verified performance. The next chapter describes the RobustIDPS.ai software platform that implements all methods as a unified, deployable system.

